\section{Results}
% Present MEASURED results. The underlying numbers come from the project's
% records (author-side), but never name internal files/fields in the paper —
% report them as observations. Pair tables with a chart for intuition.

\begin{table}[h]\centering
\begin{tabular}{lrr}
\toprule
Condition & Metric 1 & Metric 2 \\
\midrule
A & 0 & 0 \\   % <- fill from measured data
B & 0 & 0 \\
C & 0 & 0 \\
\bottomrule
\end{tabular}
\caption{Numerical comparison of the measured results.}
\end{table}

\begin{figure}[h]\centering
  \includegraphics[width=0.7\linewidth]{figures/results-bar}
  \caption{Comparison chart (edit figures/results-bar.fig.tex with real values).}
  \label{fig:chart}
\end{figure}
